
\chapter{Future Scope}

MetaGPT opens several avenues for future enhancement and research. One promising direction is to expand the set of supported project templates beyond React and Next.js to include mobile applications, microservices, and data engineering pipelines. This would require extending the agent SOPs and project scaffolding logic to accommodate additional frameworks, languages, and deployment targets.

Another important area is deeper integration with developer workflows, including Git-based version control, continuous integration (CI) pipelines, and deployment platforms. Features such as automatic commit generation, pull request creation, and environment provisioning would further reduce friction between generated projects and real-world engineering practices. Additionally, incorporating richer feedback loops---for example, using runtime metrics, static analysis tools, or human-in-the-loop annotations---could help the agents iteratively improve code quality and performance.

From a research perspective, MetaGPT can serve as a testbed for studying collaboration between multiple LLM agents, long-lived project memory, and adaptive SOPs that evolve over time. Exploring hybrid setups where human experts supervise or override specific agents, as well as experimenting with alternative LLM backends, are also valuable directions for future work.
